\section{Product Validation Strategy}

Kungfu separates evidence classes so that an impressive demonstration does
not silently become a production claim.

\subsection{A Ladder of Evidence}

\begin{evidenceladder}
  \item \textbf{Contract evidence} proves that product surfaces and authority
  objects obey checked structural rules.
  \item \textbf{Source qualification} proves that the exact source tree passes
  declared platform and subsystem gates.
  \item \textbf{Installed-artifact evidence} proves that one retained build
  artifact performs the bounded behavior under test.
  \item \textbf{Provider-bound evidence} identifies the selected executable,
  version, runtime profile, platform, plan, attempts, assessment, and report.
  \item \textbf{Longitudinal evidence} tests real Work across time, failure,
  replacement, recovery, and human intervention.
  \item \textbf{Independent adoption evidence} requires a receiver or Hub not
  controlled by the founding implementation.
\end{evidenceladder}

Higher evidence classes do not erase the need for lower ones, and lower
classes do not inherit the claims of higher ones.

\subsection{Agent Work Lab}

The exact installed-artifact Agent Work Lab path proves one bounded result:
two distinct fresh processes can preserve one Work identity, retain a partial
first result, withhold copied chat from the second process, recover the
required state, and complete the fixture through the declared oracle
\cite{agentWorkLab}.

The report exposes each continuity check rather than offering only a cinematic
success. Same-Agent and cross-Agent modes extend the same structure to
configured providers, binding provider and execution evidence while retaining
confinement and model-quality limits.

\subsection{What the Current Evidence Does Not Prove}

Current qualification does not establish:

\begin{itemize}[leftmargin=*]
  \item that a particular model is intelligent, safe, or better than another;
  \item general production continuity or every real Project shape;
  \item complete public Kungfu v4 distribution;
  \item equivalent polished behavior on every supported platform;
  \item physical power-loss durability or distributed repair;
  \item universal native-code confinement;
  \item a public KFX marketplace or independent production Agent Hubs; or
  \item a persistent self-maintaining machine subject.
\end{itemize}

These are not rhetorical disclaimers. They determine which experiment or
release gate must be completed next.

\subsection{The Real-Work Test}

The decisive product question is whether Kungfu reduces the user's recovery
and coordination burden. A valid real-work study should compare matched Work
with and without governed continuation, retain negative outcomes, and measure
how often users must restate context, repair mistaken assumptions, reconstruct
authority, or determine whether a provider process actually completed the
intended result.

Kungfu succeeds as a product only if the deeper evidence model makes ordinary
Work easier to continue and trust.
